\documentclass[UTF8,a4paper,12pt]{ctexart}

\usepackage{amsmath,amssymb,bm}
\usepackage{booktabs}
\usepackage{geometry}
\usepackage{array}
\usepackage{float}
\usepackage{graphicx}
\usepackage{longtable}

\geometry{
  left=2.5cm,
  right=2.5cm,
  top=2.5cm,
  bottom=2.5cm
}

\setlength{\parindent}{2em}
\linespread{1.25}

\begin{document}

\section{问题二：整个烘干过程的变物性热质耦合模型}

\subsection{问题分析}

问题一在预热平衡阶段采用固定热物性参数，对药材内部温度场与水分场进行了基础描述。
但在完整烘干过程中，药材含水率持续下降，其密度、比热容、导热系数和水分扩散系数均会随状态发生变化。
因此，问题二在问题一一维径向传热传质框架的基础上，引入附录3给出的变物性关系，
建立一维径向变物性非稳态热质耦合模型。

根据题意，问题二需要描述预热平衡阶段和恒温干燥阶段的整个烘干过程，
且相关经验公式统一采用附录3。
考虑到尺寸收缩已在问题四中单独讨论，本问仍将药材视为尺寸固定的圆柱体，其几何尺寸为
\begin{equation}
L=0.25~\mathrm{m},\qquad R=0.02~\mathrm{m}.
\end{equation}

圆柱长度与半径之比为
\begin{equation}
\frac{L}{R}=12.5,
\end{equation}
且圆柱两端面积与侧面积之比为
\begin{equation}
\frac{2\pi R^2}{2\pi RL}=\frac{R}{L}=0.08.
\end{equation}
因此，对于圆柱主体区域，径向传热传质效应占主导。
本文继续选取圆柱中部截面进行研究，忽略轴向梯度及端面效应，将模型简化为关于径向坐标
$r$ 和时间 $t$ 的一维非稳态问题。

\subsection{模型假设}

为保证模型具有明确的物理含义和可计算性，作如下假设：

\begin{enumerate}
  \item 药材为均质、各向同性的圆柱体，问题二中长度和半径保持不变；
  \item 药材内部温度和水分浓度仅沿径向变化，忽略轴向梯度及端面效应；
  \item 密度、比热容、导热系数和水分扩散系数采用附录3中的经验关系，并随局部状态动态更新；
  \item 对流换热系数和对流传质系数沿用问题一给定值；
  \item 附件1中的烘房温度和水分浓度作为时变外界边界条件，离散采样点之间采用分段线性插值；
  \item 本问首先建立工程型变物性基准模型，暂不计蒸发潜热、水分迁移携带显热等高阶能量耦合效应。
\end{enumerate}

\subsection{变物性关系}

根据附录3，药材密度、比热容和导热系数分别为
\begin{align}
\rho(C) &= 650+128C,\\
c_p(C) &= 1450+2736\frac{C}{C+1},\\
k(C) &= 0.21+0.38\frac{C}{C+1}.
\end{align}

其中，$C$ 表示药材干基水分浓度，单位为 $\mathrm{kg/kg}$。

水分有效扩散系数为
\begin{equation}
D(C,T_K)
=
2.4\times10^{-3}
\exp\left(-\frac{0.45}{C}\right)
\exp\left(-\frac{3850}{T_K}\right),
\end{equation}
其中
\begin{equation}
T_K=T+273.15
\end{equation}
为药材局部绝对温度，单位为 K。

由此可见，含水率的变化会通过 $\rho(C)$、$c_p(C)$ 和 $k(C)$ 影响温度场；
温度和含水率又共同决定 $D(C,T_K)$，进而影响水分迁移速度，
因此两个场之间存在双向非线性耦合。

\subsection{控制方程}

\subsubsection{温度场方程}

在忽略内部热源及高阶相变能量项的条件下，将含水率相关热物性引入非稳态导热方程，得
\begin{equation}
\boxed{
\rho(C)c_p(C)
\frac{\partial T}{\partial t}
=
\frac{1}{r}
\frac{\partial}{\partial r}
\left[
r k(C)
\frac{\partial T}{\partial r}
\right]
}
\label{eq:heat}
\end{equation}

本模型将题目给出的 $\rho(C)$、$c_p(C)$ 和 $k(C)$ 视为宏观有效热物性，
用于刻画含水率变化对药材传热能力的影响，而不进一步区分固相、水相和孔隙气相之间的微观能量输运。

\subsubsection{水分场方程}

药材内部水分迁移采用有效扩散模型描述。
在圆柱坐标系下，有
\begin{equation}
\boxed{
\frac{\partial C}{\partial t}
=
\frac{1}{r}
\frac{\partial}{\partial r}
\left[
rD(C,T_K)
\frac{\partial C}{\partial r}
\right]
}
\label{eq:mass}
\end{equation}

式~\eqref{eq:heat} 与式~\eqref{eq:mass} 的耦合关系可概括为
\begin{equation}
C
\longrightarrow
\{\rho,c_p,k\}
\longrightarrow
T,
\qquad
(T,C)
\longrightarrow
D
\longrightarrow
C.
\end{equation}

\subsection{初始条件与边界条件}

\subsubsection{初始条件}

问题二从烘干初始时刻重新计算。
初始时药材内部温度和水分浓度均匀，因此
\begin{align}
T(r,0)&=28~^\circ\mathrm{C},\\
C(r,0)&=2.55~\mathrm{kg/kg},
\qquad 0\le r\le R.
\end{align}

\subsubsection{圆柱中心边界}

由轴对称性，在 $r=0$ 处满足
\begin{align}
\left.\frac{\partial T}{\partial r}\right|_{r=0}&=0,\\
\left.\frac{\partial C}{\partial r}\right|_{r=0}&=0.
\end{align}

\subsubsection{药材表面边界}

药材表面与烘房环境之间采用第三类边界条件。

传热边界为
\begin{equation}
\boxed{
-k(C_s)
\left.
\frac{\partial T}{\partial r}
\right|_{r=R}
=
h_T\left[T_s-T_{\infty}(t)\right]
}
\label{eq:heatbc}
\end{equation}
其中
\begin{equation}
T_s=T(R,t),\qquad
h_T=25~\mathrm{W/(m^2\cdot K)}.
\end{equation}

传质边界为
\begin{equation}
\boxed{
-D(C_s,T_s)
\left.
\frac{\partial C}{\partial r}
\right|_{r=R}
=
h_m\left[C_s-C_{\infty}(t)\right]
}
\label{eq:massbc}
\end{equation}
其中
\begin{equation}
C_s=C(R,t),\qquad
h_m=8\times10^{-7}~\mathrm{m/s}.
\end{equation}

$T_{\infty}(t)$ 与 $C_{\infty}(t)$ 分别为附件1给出的烘房温度和水分浓度。
对于任意相邻采样时刻 $t_j$ 和 $t_{j+1}$，外界变量 $X_{\infty}(t)$ 采用分段线性插值：
\begin{equation}
X_{\infty}(t)
=
X_{\infty,j}
+
\frac{t-t_j}{t_{j+1}-t_j}
\left(
X_{\infty,j+1}-X_{\infty,j}
\right),
\qquad
t_j\le t\le t_{j+1}.
\end{equation}

其中 $X_{\infty}$ 可分别代表 $T_{\infty}$ 或 $C_{\infty}$。

需要说明的是，附件1中的烘房水分浓度在本文中作为表面传质过程的等效外界驱动力，
从而与药材内部干基水分浓度构成经验型 Robin 边界条件。

\subsection{预热平衡与恒温干燥阶段的统一处理}

问题二要求整个烘干过程统一采用附录3经验公式，
因此本文不在预热平衡阶段与恒温干燥阶段人为切换控制方程或物性公式。

在烘干前期，$T_{\infty}(t)$ 和 $C_{\infty}(t)$ 随时间显著变化，
药材内部温度与含水率随之响应；
进入恒温干燥阶段后，烘房温湿环境逐渐趋于稳定，
但药材内部水分仍持续迁移，因此 $\rho(C)$、$c_p(C)$、$k(C)$ 和 $D(C,T_K)$ 仍会继续变化。

由此，两个阶段的差异通过外部边界条件和局部物性的连续变化自然体现，
无需人为设置参数突变。

\subsection{尺度分析与模型合理性}

初始状态为
\begin{equation}
C_0=2.55,\qquad T_0=28~^\circ\mathrm{C}.
\end{equation}
代入附录3得
\begin{align}
\rho_0&=976.4~\mathrm{kg/m^3},\\
c_{p,0}&\approx3415.30~\mathrm{J/(kg\cdot K)},\\
k_0&\approx0.48296~\mathrm{W/(m\cdot K)},\\
D_0&\approx5.6417\times10^{-9}~\mathrm{m^2/s}.
\end{align}

初始热扩散率为
\begin{equation}
\alpha_0
=
\frac{k_0}{\rho_0c_{p,0}}
\approx1.4483\times10^{-7}~\mathrm{m^2/s}.
\end{equation}

因此径向热扩散特征时间约为
\begin{equation}
\tau_T
\sim
\frac{R^2}{\alpha_0}
\approx0.767~\mathrm{h},
\end{equation}
而水分扩散特征时间约为
\begin{equation}
\tau_C
\sim
\frac{R^2}{D_0}
\approx19.695~\mathrm{h}.
\end{equation}

故有
\begin{equation}
\tau_T\ll\tau_C,
\end{equation}
说明温度场达到近似平衡的速度明显快于水分扩散过程，
与题目所描述的预热平衡和长期恒温干燥阶段相一致。

进一步有
\begin{equation}
Bi_T=\frac{h_TR}{k_0}\approx1.035,
\end{equation}
以及等效传质 Biot 数
\begin{equation}
Bi_C=\frac{h_mR}{D_0}\approx2.836.
\end{equation}

二者均明显大于 $0.1$，说明药材内部存在不可忽略的径向梯度，
因此不宜采用空间均匀的集中参数模型，而应保留径向偏微分方程。

\subsection{有限体积离散}

\subsubsection{空间网格}

将 $0\le r\le R$ 划分为 $N$ 个均匀控制体，
网格宽度为
\begin{equation}
\Delta r=\frac{R}{N}.
\end{equation}

第 $i$ 个控制体中心记为 $r_i$，
左右界面分别为 $r_{i-\frac12}$ 和 $r_{i+\frac12}$。
其体积与界面面积分别为
\begin{align}
V_i
&=
\pi L
\left(
r_{i+\frac12}^2-r_{i-\frac12}^2
\right),\\
A_{i\pm\frac12}
&=
2\pi Lr_{i\pm\frac12}.
\end{align}

\subsubsection{温度方程离散}

采用后向欧拉全隐式格式。
在时间层 $t^{n+1}=t^n+\Delta t$，对式~\eqref{eq:heat} 在控制体内积分，可得
\begin{equation}
\rho_i c_{p,i}V_i
\frac{T_i^{n+1}-T_i^n}{\Delta t}
=
G_{i-\frac12}^{T}
\left(T_{i-1}^{n+1}-T_i^{n+1}\right)
+
G_{i+\frac12}^{T}
\left(T_{i+1}^{n+1}-T_i^{n+1}\right),
\label{eq:heatfvm}
\end{equation}
其中内部界面导热系数采用调和平均
\begin{equation}
k_{i+\frac12}
=
\frac{2k_ik_{i+1}}{k_i+k_{i+1}},
\end{equation}
对应热传导系数为
\begin{equation}
G_{i+\frac12}^{T}
=
\frac{k_{i+\frac12}A_{i+\frac12}}{\Delta r}.
\end{equation}

这种处理能够保证相邻控制体界面热流连续，并减小变物性条件下简单算术平均造成的误差。

\subsubsection{水分方程离散}

对式~\eqref{eq:mass} 采用相同的后向欧拉有限体积格式：
\begin{equation}
V_i
\frac{C_i^{n+1}-C_i^n}{\Delta t}
=
G_{i-\frac12}^{C}
\left(C_{i-1}^{n+1}-C_i^{n+1}\right)
+
G_{i+\frac12}^{C}
\left(C_{i+1}^{n+1}-C_i^{n+1}\right),
\label{eq:massfvm}
\end{equation}
其中
\begin{equation}
D_{i+\frac12}
=
\frac{2D_iD_{i+1}}{D_i+D_{i+1}},
\end{equation}
以及
\begin{equation}
G_{i+\frac12}^{C}
=
\frac{D_{i+\frac12}A_{i+\frac12}}{\Delta r}.
\end{equation}

由于 $D=D(C,T_K)$，水分方程属于非线性变系数扩散方程。

\subsubsection{表面 Robin 边界离散}

对于最外层控制体，将半个控制体内部传递阻力与外部对流阻力串联。

温度边界的等效传热系数为
\begin{equation}
G_R^T
=
\frac{A_R}
{\dfrac{\Delta r}{2k_s}+\dfrac{1}{h_T}},
\end{equation}
因此最外层温度方程中的边界贡献为
\begin{equation}
G_R^T
\left(
T_{\infty}^{n+1}-T_N^{n+1}
\right).
\end{equation}

药材表面温度由
\begin{equation}
q_R^T
=
\frac{T_N^{n+1}-T_{\infty}^{n+1}}
{\dfrac{\Delta r}{2k_s}+\dfrac{1}{h_T}}
\end{equation}
反算为
\begin{equation}
T_s^{n+1}
=
T_{\infty}^{n+1}
+
\frac{q_R^T}{h_T}.
\end{equation}

类似地，水分边界的等效传质系数为
\begin{equation}
G_R^C
=
\frac{A_R}
{\dfrac{\Delta r}{2D_s}+\dfrac{1}{h_m}},
\end{equation}
边界贡献为
\begin{equation}
G_R^C
\left(
C_{\infty}^{n+1}-C_N^{n+1}
\right).
\end{equation}

定义等效水分通量
\begin{equation}
j_R
=
\frac{C_N^{n+1}-C_{\infty}^{n+1}}
{\dfrac{\Delta r}{2D_s}+\dfrac{1}{h_m}},
\end{equation}
则表面水分浓度为
\begin{equation}
C_s^{n+1}
=
C_{\infty}^{n+1}
+
\frac{j_R}{h_m}.
\end{equation}

此处的 $j_R$ 是以干基水分浓度为变量构造的等效通量，
不将其直接解释为真实水质量通量。

\subsection{非线性耦合 Picard 迭代}

由于 $\rho$、$c_p$、$k$ 依赖 $C$，而 $D$ 同时依赖 $C$ 和 $T$，
每个时间层均需要进行非线性耦合迭代。

在时间层 $n+1$，首先令
\begin{equation}
T^{(0)}=T^n,\qquad C^{(0)}=C^n.
\end{equation}

第 $m$ 次迭代按以下顺序进行：

\begin{enumerate}
  \item 根据 $C^{(m)}$ 更新 $\rho^{(m)}$、$c_p^{(m)}$ 和 $k^{(m)}$；
  \item 组装温度有限体积方程，求得 $T^{(m+1)}$；
  \item 根据 $C^{(m)}$ 与新得到的 $T^{(m+1)}$ 更新
  \[
  D^{(m)}=D\left(C^{(m)},T^{(m+1)}\right);
  \]
  \item 组装水分有限体积方程，求得 $C^{(m+1)}$；
  \item 更新表面 $T_s$、$C_s$ 以及相应的边界物性；
  \item 检查收敛条件。
\end{enumerate}

本文采用无穷范数收敛准则
\begin{equation}
\max
\left\{
\left\|T^{(m+1)}-T^{(m)}\right\|_{\infty},
\left\|C^{(m+1)}-C^{(m)}\right\|_{\infty},
|T_s^{(m+1)}-T_s^{(m)}|,
|C_s^{(m+1)}-C_s^{(m)}|
\right\}
<\varepsilon,
\end{equation}
并取
\begin{equation}
\varepsilon=10^{-10}.
\end{equation}

每次迭代得到的温度和水分方程均为三对角线性方程组，
采用追赶法求解。

\subsection{网格、时间步长及输出重构}

最终基准计算采用
\begin{equation}
N=400,\qquad
\Delta r=5.0\times10^{-5}~\mathrm{m},
\qquad
\Delta t=0.0625~\mathrm{s}.
\end{equation}

计算过程中最大 Picard 迭代次数为 4 次，所有时间步均正常收敛。

由于有限体积未知量位于控制体中心，
题目要求的 $r=0$ 和 $r=R$ 并非直接未知量。
轴心值利用轴对称条件进行局部二次重构：
\begin{equation}
u(0)
=
\frac{9u_1-u_2}{8},
\end{equation}
其中 $u$ 可代表 $T$ 或 $C$。

内部指定位置采用相邻控制体中心线性插值，
表面值则由 Robin 边界关系直接反算。

\subsection{数值收敛性检验}

为验证数值结果对空间网格和时间步长的稳定性，
在题目要求的 6 个时刻和 5 个径向位置上比较不同离散尺度的结果，
以所有输出点的最大绝对差作为误差指标。

\begin{table}[H]
\centering
\caption{问题二数值收敛性检验}
\begin{tabular}{lccc}
\toprule
检验项目 & 对比设置 & 最大温度差/($^\circ$C) & 最大水分差/(kg/kg)\\
\midrule
空间网格细化
& $N=400\rightarrow800,\ \Delta t=0.0625$ s
& $6.66\times10^{-6}$
& $5.00\times10^{-6}$\\
时间步长细化
& $\Delta t=0.125\rightarrow0.0625$ s，$N=400$
& $9.23\times10^{-5}$
& $5.71\times10^{-6}$\\
\bottomrule
\end{tabular}
\end{table}

上述差异均小于题目四位小数输出精度对应的量级，
说明 $N=400$、$\Delta t=0.0625$ s 已能满足问题二的数值精度要求。

\subsection{计算结果}

\subsubsection{3 h 内药材温度分布}

采用上述模型和数值设置计算得到表~\ref{tab:q2T}。

\begin{table}[H]
\centering
\caption{3 h 内药材温度}
\label{tab:q2T}
\begin{tabular}{cccccc}
\toprule
时间/h & 0 cm & 0.5 cm & 1 cm & 1.5 cm & 2 cm\\
\midrule
0.5 & 32.1893 & 32.3822 & 32.9660 & 33.9609 & 35.4131\\
1.0 & 40.3816 & 40.5536 & 41.0605 & 41.8782 & 42.9976\\
1.5 & 45.8467 & 45.9348 & 46.1929 & 46.6050 & 47.1400\\
2.0 & 48.4502 & 48.4880 & 48.5984 & 48.7735 & 49.0033\\
2.5 & 49.4670 & 49.4792 & 49.5136 & 49.5653 & 49.6609\\
3.0 & 49.8495 & 49.8553 & 49.8746 & 49.9101 & 49.9664\\
\bottomrule
\end{tabular}
\end{table}

结果表明，烘干初期药材表面温度高于中心温度，
径向温度梯度较明显；
随着时间增加，药材整体温度逐渐接近烘房约 $50^\circ\mathrm{C}$ 的稳定水平。
到 3 h 时，中心温度约为 $49.8495^\circ\mathrm{C}$，
表面温度约为 $49.9664^\circ\mathrm{C}$，
径向温差已明显减小，表明温度场已基本进入近似平衡状态。

\subsubsection{3 h 内药材水分浓度分布}

计算得到的水分浓度见表~\ref{tab:q2C}。

\begin{table}[H]
\centering
\caption{3 h 内药材水分浓度}
\label{tab:q2C}
\begin{tabular}{cccccc}
\toprule
时间/h & 0 cm & 0.5 cm & 1 cm & 1.5 cm & 2 cm\\
\midrule
0.5 & 2.5499 & 2.5489 & 2.5256 & 2.3257 & 1.6486\\
1.0 & 2.5257 & 2.4948 & 2.3578 & 2.0230 & 1.4711\\
1.5 & 2.3861 & 2.3256 & 2.1344 & 1.8020 & 1.3476\\
2.0 & 2.1709 & 2.1084 & 1.9236 & 1.6259 & 1.2311\\
2.5 & 1.9566 & 1.9006 & 1.7360 & 1.4720 & 1.1166\\
3.0 & 1.7662 & 1.7165 & 1.5702 & 1.3333 & 1.0081\\
\bottomrule
\end{tabular}
\end{table}

与温度场相比，水分场在 3 h 时仍具有明显的径向梯度。
表面水分浓度已下降至约 $1.0081~\mathrm{kg/kg}$，
而中心仍约为 $1.7662~\mathrm{kg/kg}$。
这说明水分迁移的时间尺度显著长于热扩散时间尺度，
与前述 $\tau_T\ll\tau_C$ 的尺度分析结果一致。

\subsection{模型讨论与局限性}

本问建立的模型利用题目给出的经验物性关系描述药材宏观传热传质行为，
并在固定几何条件下实现了温度场和水分场的双向变物性耦合。

需要指出的是，题给
\begin{equation}
\rho=650+128C
\end{equation}
与固定体积、干物质密度严格不变的多组分质量关系并不完全兼容，
因此本文将其视为宏观有效物性关系，
不进一步将其解释为所有组分均严格闭合的微观密度模型。

同时，本基准模型暂未引入蒸发潜热。
已有干燥研究表明，潜热在部分物料中可能显著影响温度预测，
但若要在本题中加入该项，需要首先将当前以干基水分浓度为变量的等效传质通量
一致转换为真实水质量通量。
在缺乏额外干物质体积密度或气固平衡关系的条件下，
直接加入潜热反而可能引入新的参数不确定性。
因此，本问先采用上述工程型变物性模型作为基准模型，
后续可通过敏感性分析进一步判断潜热修正的必要性。

\end{document}
